\section{Identification of the limiting field}
\label{sec:limit}

Fix \(x\) throughout, and let \(v\in(0,\infty)\) be the
domain-independent constant in Proposition~\ref{prop:plane-covariance}.
We identify every subsequential limit of the actual zero-boundary
fields.  The steps below also establish tightness and justify passage
of all moments; no assertion of Gaussianity is used before the last
step.

\subsection{Compactness, moments, and the insertion identity}

We write \(h_\delta(f)=\int_D h_\delta(z)f(z)\dd z\).
This is the exact face-area pairing from the introduction, including
for signed tests.  In particular \(\int f=0\) means that its discrete
weights sum to exactly zero, so adding a constant to a height lift has
no effect.

\begin{proposition}[Tightness and moment convergence]
\label{prop:tightness}
For every \(t>1\), the true fields are tight in
\(H^{-t}_{\mathrm{loc}}(D)\).  The same holds for the upper and lower
brackets, uniformly in their exceptional windows.  Every fixed finite
collection of test pairings has uniformly bounded moments of every
order.  Along a distributional law limit, all its mixed moments
converge.  These assertions hold also for finite joint couplings of
the indicated fields.
\end{proposition}

\begin{proof}
A consequence of Proposition~\ref{prop:smearing}, using a finite
cover of a compact set by interior disks, is
\begin{equation}
 \sup_\delta\|h_\delta(f)\|_{L^p}
 \le C_{K,p}\|f\|_\infty,\qquad
 \supp f\subset K\Subset D,\quad 1\le p<\infty.
 \label{eq:limit-fixed-moments}
\end{equation}
Indeed, split the face weights by a smooth partition of unity
subordinate to that cover.  Each resulting weight family has bounded
total variation and the required maximum weight on its fixed disk;
apply the smearing estimate and the triangle inequality.  The argument
applies to the brackets after conditioning on their traces, whose
values lie between \(-1\) and \(1\).  Its constants thus do not depend
on the windows.  Complex tests are handled by their real and imaginary
parts.

Here is the compactness argument in detail.  Place \(D\) in a square
torus and fix \(\chi\in C_c^\infty(D)\).  For its Fourier basis
\(e_k\), \(k\in\Z^2\), let
\(a_{\delta,k}=h_\delta(\chi\overline{e_k})\).  Since the moduli of the
basis functions are constant, \eqref{eq:limit-fixed-moments} bounds
\(\|a_{\delta,k}\|_{L^p}\) independently of \(k\).  For \(s>1\) and
\(p\ge2\), Minkowski's inequality gives
\begin{align}
 \bigl\|\|\chi h_\delta\|_{H^{-s}}^2\bigr\|_{L^{p/2}}
 &\le
 \sum_{k\in\Z^2}(1+|k|^2)^{-s}
             \|a_{\delta,k}\|_{L^p}^2
 \le C_{\chi,p,s}.
 \label{eq:limit-sobolev-moments}
\end{align}
The bounds for smaller \(p\) follow by monotonicity of \(L^p\) norms.
Choose \(1<s<t\).  A bounded \(H^{-s}\) ball is relatively compact
in \(H^{-t}\): its Fourier tail beyond frequency \(M\) has squared
\(H^{-t}\) norm at most
\((1+M^2)^{-(t-s)}\) times its squared \(H^{-s}\) norm, and its
first finitely many coordinates lie in a bounded finite-dimensional
set.  Markov's inequality and \eqref{eq:limit-sobolev-moments}
therefore give tightness of \(\chi h_\delta\) in \(H^{-t}\).

Choose nested interior compacts exhausting \(D\), and smooth cutoffs
equal to one on successive compacts.  For any prescribed error
probability, choose a compact set for the \(n\)-th cutoff with error
at most that probability times \(2^{-n}\).  The simultaneous
localizations then lie in a relatively compact set for the local
\(H^{-t}\) topology.  This proves the stated local tightness.
Equivalently, this topology is the compatible subspace of a countable
product of the localized Sobolev spaces, with its usual metrizable
topology.  The same argument and a union bound treat finite joint
couplings.

For fixed \(f_1,\ldots,f_k\), a monomial
\(\prod_i h_\delta(f_i)^{a_i}\) has uniformly bounded \(L^{1+\epsilon}\)
norm, for instance with \(\epsilon=1\), by
\eqref{eq:limit-fixed-moments} and H\"older's inequality at sufficiently
high orders.  Its uniform integrability, joint convergence of the
pairings, and truncation prove convergence of that mixed moment.
The same argument gives moment convergence with bounded continuous
functions of finitely many pairings.  In particular all the preceding
moment bounds pass to every limit.
\end{proof}

A second consequence of the smearing estimate, which we will use at
collisions, is the following.  If
\(f_{z,r}(w)=r^{-2}f((w-z)/r)\), where \(f\) ranges over a fixed bounded
family of smooth functions supported in a fixed disk, then on an
interior compact, for sufficiently small \(r\),
\begin{equation}
 \|h_\delta(f_{z,r})\|_{L^p}
 \le C_{p,f}\bigl(1+\log(1/r)\bigr)^{1/2}.
 \label{eq:limit-small-moments}
\end{equation}
It suffices here and below to take the mesh sufficiently small at each
fixed \(r\).  Proposition~\ref{prop:smearing} gives this bound since
the distance to the boundary is bounded below and the smearing radius
is \(r\).  It passes to limits by Proposition~\ref{prop:tightness}.
For plane fields and zero-total \(f\), Proposition~\ref{prop:plane-moments}
instead gives a bound independent of \(r\).

Fix henceforth any subsequence on which the true fields converge,
and denote its limit by \(H\).  Height reflection and moment
convergence imply \(\E H(f)=0\).

\begin{lemma}[Transfer to the chosen field]
\label{limit:transfer}
The field \(H\) satisfies the limiting insertion identity of
Proposition~\ref{prop:bracket-insertion} in every disk arrangement
allowed there:
\begin{equation}
 \E\left[H(\Delta\varphi)
          \prod_{\ell=1}^k H(g_\ell)\right]=0,
 \qquad \int g_\ell=0,
 \label{eq:limit-insertion}
\end{equation}
where each \(g_\ell\) is supported in its indicated spectator disk.
The same assertion first holds with each spectator replaced by a
bounded continuous function of finitely many zero-total tests in that
disk.
\end{lemma}

\begin{proof}
Fix a finite grid of regular horizontal levels and positive exceptional
windows sufficiently small for the given arrangement.  We pass first
to the mesh limit and then to the window limit, keeping the chosen
true-field marginal throughout.

Couple its lower and upper brackets with the true field in order,
as in Lemma~\ref{pins:bracket-order}.  First take a joint law limit
along the chosen subsequence, with the windows fixed.  Next shrink
the windows and take a further joint limit.  Proposition~\ref{prop:tightness}
applies uniformly throughout, so these extractions are possible and
the middle marginal is always the already chosen law of \(H\).
Distributional order passes to the limits because evaluation against
each nonnegative smooth test is continuous.  Write the resulting
ordered triple as
\[
       H^-\le H\le H^+.
\]
Proposition~\ref{prop:bracket-collapse} and moment convergence give
\(\E H^-(f)=\E H^+(f)=0\), as well as \(\E H(f)=0\).
For every nonnegative smooth \(f\), both nonnegative random variables
\(H^+(f)-H(f)\) and \(H(f)-H^-(f)\) thus have expectation zero.
Apply this to a countable dense family of nonnegative tests on an
exhaustion of \(D\).  Continuity of distributions then shows
\(H^-=H=H^+\) almost surely as distributions.

At fixed windows, Proposition~\ref{prop:bracket-insertion} gives the
insertion identity for bounded local increment factors.  For bounded
continuous functions of zero-total pairings it passes to the first
law limit by uniform integrability of \(H(\Delta\varphi)\), and then
to the window limit by the same argument.  The collapsed marginal is
the chosen \(H\), so this proves that identity for \(H\).
Repeat for any desired finite grid; no change to the chosen marginal
is involved.  Finally truncate each \(H(g_\ell)\) by a bounded
continuous truncation.  H\"older's inequality and moments of orders
larger than the number of factors show that the expectation of the
truncation error tends to zero.  This proves
\eqref{eq:limit-insertion}.
\end{proof}

\subsection{Boundary references and harmonic moment distributions}

The insertion identity currently permits only zero-total spectator
tests.  To obtain harmonicity of the full moment distributions, we
must allow arbitrary tests, including averages of absolute heights.
We construct reference averages that vanish near the boundary;
subtracting such a reference expresses each spectator as a finite
sum of local zero-total observations.

Fix a nonnegative smooth radial bump \(\rho\) of integral one,
supported in the unit disk.  All unit radial averages below are
translations and dilations of this bump.

\begin{lemma}[Vanishing references]
\label{limit:references}
There are finite averages \(R_m\) of unit radial pairings of \(H\),
supported arbitrarily near the boundary, such that
\[
       \|R_m\|_{L^p}\longrightarrow0
       \quad\text{for every }p<\infty.
\]
Each \(R_m\) uses finitely many distinct boundary points, with depth
chosen after those points are fixed.  Its constituent disks can be
required to have any prescribed small positive maximum diameter.
\end{lemma}

\begin{proof}
Choose distinct points \(p_1,\ldots,p_m\) on a fixed boundary arc.
For small \(a>0\), put the centers a distance \(a\) along their inward
normals and use bumps \(\psi_{\ell,a}\) of radius \(\theta a\), where
\(0<\theta<1/4\) is fixed.  The tubular neighborhood theorem for a
\(C^2\) Jordan curve gives distance exactly \(a\) to the boundary for
these centers when \(a\) is sufficiently small.  Define
\[
 Y_{\ell,a}=H(\psi_{\ell,a}),\qquad
 R_{m,a}=\frac1m\sum_{\ell=1}^m Y_{\ell,a}.
\]
The comparable-depth moment bound in
Proposition~\ref{prop:boundary-moments} is uniform along this fixed
arc.  Minkowski's inequality therefore gives
\(\sup_{m,a}\|R_{m,a}\|_{L^q}\le C_q\) for each finite \(q\), after
restricting to the corresponding sufficiently small depths.
For fixed \(m\), the same proposition gives
\(\E[Y_{\ell,a}Y_{k,a}]\to0\) when \(\ell\ne k\).
Consequently
\begin{equation}
  \limsup_{a\downarrow0}\|R_{m,a}\|_{L^2}^2
       \le \frac{C_2^2}{m}.
  \label{eq:limit-reference-second}
\end{equation}
For each \(m\), choose a positive depth \(a_m\) small enough that this
squared norm is at most \((C_2^2+1)/m\) and all disk diameters meet
the prescribed bound.  Then \(R_m=R_{m,a_m}\) tends to zero in
\(L^2\).  For \(p>2\), interpolate this convergence with the uniform
\(L^q\) bound for any \(q>p\); for \(p\le2\), use monotonicity.
This proves the assertion.  The order in
\eqref{eq:limit-reference-second} is depth first at fixed \(m\),
then \(m\to\infty\).
\end{proof}

For \(j\ge1\), define the moment distribution \(T_j=\E H^{\otimes j}\)
on \(D^j\), and set \(T_0=1\).  This definition is legitimate for
general smooth tests on the product, not just tensor tests.
For example, after localizing in a square torus, expand a product-space
test in its Fourier series.  Its coefficients decay faster than any
power, while H\"older's inequality and
\eqref{eq:limit-fixed-moments} bound the expected absolute product of
any \(j\) Fourier coefficients uniformly.  The resulting series is
absolutely convergent and defines a continuous distribution.
For tensor tests it gives precisely
\[
 T_j(f_1\otimes\cdots\otimes f_j)
       =\E\prod_{i=1}^j H(f_i).
\]
This construction agrees with taking the expectation of the tensor
power of the random distribution by
\eqref{eq:limit-sobolev-moments}.

For \(j\ge2\), we will determine \(T_j\) from \(T_{j-2}\) by its behavior
in one variable.  The next steps establish harmonicity away from coincident
points, rule out additional distributions supported at coincidences,
and identify each logarithmic collision coefficient.  Together with
the zero boundary value, these properties will give the Gaussian
moment recursion by subtraction of Dirichlet Green functions.

\begin{proposition}[Harmonic moments]
\label{prop:harmonic-moments}
Let
\[
 \mathcal C_j=\{(z_1,\ldots,z_j)\in D^j:
                         z_i\ne z_k\text{ for }i\ne k\}.
\]
On \(\mathcal C_j\), \(T_j\) is a smooth function and is harmonic in
each variable separately.
\end{proposition}

\begin{proof}
We prove the first-variable distributional Laplace identity in a
small product neighborhood of any distinct tuple; symmetry then gives
all the identities.  Translate the first center to the origin.
Choose \(R>0\) so that \(B(0,4R)\Subset D\) and all the other centers
lie outside \(\overline{B(0,4R)}\).  Choose \(0<b<R/1000\) and
\(0<\eta<R/(1000j)\).  Fix the finite regular grid supplied by
Lemma~\ref{pins:local-geometry} for \(z_0=0\), \(\rho=R\), and
\(q=j-1\).  Its choice is independent of the small spectator disks
and chains constructed below.
Take \(\varphi\) supported in \(B(0,b)\), and take each spectator
test \(f_i\), \(2\le i\le j\), in a sufficiently small neighborhood
of its center, of diameter less than \(\eta/4\).
We will show
\begin{equation}
       \E\left[H(\Delta\varphi)\prod_{i=2}^j H(f_i)\right]=0.
       \label{eq:limit-harmonic-tensor}
\end{equation}
For \(j=1\) this follows from the zero mean; suppose \(j\ge2\).

Write \(c_i=\int f_i\), and choose a unit bump \(\rho_i\) in that
spectator neighborhood.  Take a reference \(R_{m,a}\) from the
preceding lemma, at a fixed positive depth small enough that all
its disks have diameter less than \(\eta/4\) and lie outside
\(\overline{B(0,4R)}\).  This is possible because that closed ball
is a compact subset of \(D\), whereas the reference disks can be
placed arbitrarily near the boundary.
Every center can be joined to each reference center by a finite path
in \(D\setminus\overline{B(0,R)}\).  To verify this last geometric
point, start with a polygonal path in the connected open set \(D\).
Whenever it enters the protected disk, replace that portion by an
arc in \(B(0,2R)\setminus\overline{B(0,R)}\); this annulus lies in
\(D\).  A small perturbation and a finite subdivision give a path
with compact image and positive clearance from both the protected
disk and the boundary.

Along each such path choose a finite chain of smooth unit bumps,
beginning with \(\rho_i\) and ending with \(\psi_{\ell,a}\).
Make the step lengths and bump radii so small that the union of
two consecutive supports is contained in a disk of radius less
than \(\eta\), disjoint from \(\overline{B(0,R)}\).
If necessary interpolate their radii in finitely many steps near
the endpoints.  Telescoping yields
\[
 H(f_i)-c_iH(\psi_{\ell,a})
  =H(f_i-c_i\rho_i)
    +c_i\sum_{\text{chain}}H(\rho_{\mathrm{old}}-\rho_{\mathrm{new}}).
\]
Every summand on the right is a zero-total test in one of these
small disks.  Averaging this equality over \(\ell\) expresses
\(H(f_i)-c_iR_{m,a}\) as a finite linear combination of such
observations.

Expand the product of these \(j-1\) finite sums.  Each monomial
selects at most \(j-1\) spectator disks, whatever the number or
length of the chains.  The insertion support lies in
\(B(0,R/1000)\), and every selected disk is outside \(B(0,R)\)
with radius less than \(R/(1000j)\).  These are precisely the
hypotheses for the grid already fixed from
Lemma~\ref{pins:local-geometry}.  For each monomial its two selected
levels give the oblique separation of the middle-strip boundary
pieces and spectators needed by Proposition~\ref{prop:bracket-insertion}.
The pair of levels may depend on the monomial; the grid does not.

There are only finitely many monomials.  Choose the exceptional windows
small enough to preserve all their positive margins.
Lemma~\ref{limit:transfer} gives zero expectation for each monomial
with the insertion.  Summing these exact identities gives
\begin{equation}
 \E\left[H(\Delta\varphi)
              \prod_{i=2}^j\bigl(H(f_i)-c_iR_{m,a}\bigr)\right]=0.
 \label{eq:limit-subtracted-references}
\end{equation}
All disks, chains, tests and grid levels used to obtain this equality
are finite and fixed before using the limiting insertion identity.

It remains to remove the reference.  This error is estimated before
the chain expansion.  The algebraic difference between the product
in \eqref{eq:limit-subtracted-references} and \(\prod_iH(f_i)\)
is a sum of \(j-1\) terms, each containing one factor \(c_iR_{m,a}\)
and other factors of the form \(H(f_k)\) or
\(H(f_k)-c_kR_{m,a}\).  H\"older's inequality and the uniform
moment bounds give, for some finite \(p\) depending only on \(j\),
\[
 \left|\E H(\Delta\varphi)
       \left\{\prod_{i=2}^jH(f_i)
        -\prod_{i=2}^j(H(f_i)-c_iR_{m,a})\right\}\right|
       \le C\|R_{m,a}\|_{L^p}.
\]
The constant does not count chain terms.  Choose depths and then
\(m\) as in Lemma~\ref{limit:references}; the right side tends to
zero.  This proves \eqref{eq:limit-harmonic-tensor}.

Finite sums of tensor tests are dense in the test-function space
on each such product neighborhood, for example by a localized
Fourier expansion with convergence of every derivative.
Consequently \(\Delta_{z_1}T_j=0\) there as a distribution.
Symmetry gives all \(j\) separate Laplace equations.  Their sum is
the ordinary elliptic Laplacian on \(\R^{2j}\); the local regularity
theorem for harmonic distributions makes \(T_j\) smooth there.
The separate distributional equations then hold pointwise.
\end{proof}

\begin{lemma}[No distributions supported on collisions]
\label{lem:no-diagonal}
On every compact subset of \(D^j\), the off-diagonal function from
Proposition~\ref{prop:harmonic-moments} is locally integrable and
represents all of \(T_j\).  For \(j\ge2\), if
\(d(\mathbf z)=\min_{i<k}|z_i-z_k|\), then, locally away from
the boundary,
\begin{equation}
 |T_j(\mathbf z)|
 \le C\bigl(1+\log_+(1/d(\mathbf z))\bigr)^{j/2},
 \qquad \mathbf z\in\mathcal C_j.
 \label{eq:limit-log-bound}
\end{equation}
Moreover, for distinct fixed \(z_2,\ldots,z_j\in D\),
\begin{equation}
 T_j(z_1,z_2,\ldots,z_j)\longrightarrow0
        \quad\text{as }z_1\longrightarrow p\in\partial D.
 \label{eq:limit-zero-boundary}
\end{equation}
\end{lemma}

\begin{proof}
Let \(T_{j,h}\) be the convolution of \(T_j\) with a product of
\(j\) radial unit bumps of common radius \(h\), considered on a
fixed interior compact and with \(h\) small enough to remain in
\(D\).  By H\"older's inequality and
\eqref{eq:limit-small-moments},
\begin{equation}
 |T_{j,h}(\mathbf z)|
   =\left|\E\prod_{i=1}^jH(\rho_{z_i,h})\right|
   \le C(1+\log(1/h))^{j/2},
 \label{eq:limit-mollifier-bound}
\end{equation}
including at coincident centers.  Whenever \(d(\mathbf z)>2h\),
the product of bump supports is off the diagonals.  The radial mean
value property, successively in each harmonic variable, gives
\(T_{j,h}(\mathbf z)=T_j(\mathbf z)\).
Taking \(h\) to be the smaller of \(d(\mathbf z)/4\) and a fixed
interior radius proves \eqref{eq:limit-log-bound}.

For completeness, a collision tube
\[
       A_h=\{\mathbf z:d(\mathbf z)\le2h\}
\]
in a fixed compact set has volume at most \(C_jh^2\): use the
union bound over pairs, fix all but one point in a selected pair,
and integrate that point over a disk of radius \(2h\).
The same bound at every radius \(s\), followed by integration of
the distribution function, gives for \(q\ge0\)
\[
 \int_{A_h}\bigl(1+\log_+(1/d(\mathbf z))\bigr)^q\dd\mathbf z
       \le C_qh^2(1+\log(1/h))^q.
\]
Indeed for \(q>0\), integrate
\(q(1+\log(1/s))^{q-1}\dd s/s\) against the bound \(Cs^2\);
the endpoint term is \(Ch^2(1+\log(1/h))^q\).
This proves local integrability.  It also shows that the
off-diagonal function has vanishing \(L^1\) mass on \(A_h\).
Equation~\eqref{eq:limit-mollifier-bound} gives the same assertion
for \(T_{j,h}\), since \(h^2(1+\log(1/h))^{j/2}\to0\).
Off \(A_h\) the two functions agree exactly.  Thus \(T_{j,h}\)
converges locally in \(L^1\) to the off-diagonal function.
But convolution also converges to \(T_j\) in distributions.
These limits are equal.  In particular neither a delta distribution
nor any of its derivatives on a collision diagonal can be present.
The case \(j=1\) has no collision set and follows directly from
harmonicity.

To prove \eqref{eq:limit-zero-boundary}, choose disjoint small fixed
radial averaging disks around \(z_2,\ldots,z_j\).
For \(z_1\) near \(p\), choose its averaging radius to be
\(\theta\dist(z_1,\partial D)\).
The supports are pairwise separated.  The harmonic mean value
property in every variable identifies the displayed pointwise
value with
\[
 \E\left[H(\rho_{z_1,\theta\dist(z_1,\partial D)})
                         \prod_{i=2}^jH(\rho_{z_i,s_i})\right].
\]
Proposition~\ref{prop:boundary-moments} makes this tend to zero.
For \(j=1\) the same conclusion also follows from the zero mean.
\end{proof}

\subsection{The logarithmic collision coefficient}

\begin{lemma}[Punctured harmonic expansion]
\label{limit:puncture}
Let \(j\ge2\).  Fix a collision point \(y\in D\) and distinct spectators
\(\mathbf z=(z_3,\ldots,z_j)\), all separated from \(y\).
For \(x\) in a sufficiently small punctured disk about \(y\),
\begin{equation}
 T_j(x,y,\mathbf z)
      =L(y,\mathbf z)\log(1/|x-y|)
                           +R(x;y,\mathbf z),
 \label{eq:limit-puncture}
\end{equation}
where \(R\) extends harmonically across \(x=y\).
The coefficient \(L\) is smooth in the separated parameters
\((y,\mathbf z)\), and the first-variable derivatives of \(R\)
are bounded uniformly on smaller disks and compact sets of these
parameters.  The coefficient \(L\) is separately harmonic in each
spectator variable.
\end{lemma}

\begin{proof}
For fixed parameters, the angular Fourier coefficients of a harmonic
function in a punctured disk have the forms
\(a_0+b_0\log r\) and \(a_nr^n+b_nr^{-n}\), \(n\ge1\), for the
sine and cosine modes.  Equation~\eqref{eq:limit-log-bound} bounds
the function by a power of \(1+\log(1/r)\).  Multiplying its
\(n\)-th Fourier coefficient by \(r^n\) and sending \(r\downarrow0\)
therefore shows \(b_n=0\) for every \(n\ge1\).
After subtracting the logarithmic term, the remaining nonnegative
modes are the harmonic extension of their values on any fixed
smaller circle, and hence extend across the center.

On a circle of radius \(s>0\) enclosing no other spectator,
the coefficient is the flux
\begin{equation}
 L(y,\mathbf z)
   =-\frac1{2\pi}
       \int_{|x-y|=s}\partial_{n_x}T_j(x,y,\mathbf z)\dd\ell(x).
 \label{eq:limit-flux}
\end{equation}
One can keep \(s\) fixed on a neighborhood of any separated
parameter tuple.  Parameterizing \(x=y+se^{i\theta}\), the integrand
is smooth by Proposition~\ref{prop:harmonic-moments}; thus \(L\)
is smooth.  Applying a spectator Laplacian under this integral
also proves its separate harmonicity.
Subtract \(L\log(1/|x-y|)\) on that circle.  Its values are
smooth and uniformly bounded on compact parameter sets.
The Poisson formula for their harmonic extension, differentiated
in \(x\) on a smaller disk, yields the claimed uniform derivative
bounds for \(R\).
\end{proof}

\begin{proposition}[Collision coefficient]
\label{prop:collision}
In \eqref{eq:limit-puncture},
\begin{equation}
              L(y,\mathbf z)=\frac{v}{2\pi}T_{j-2}(\mathbf z).
 \label{eq:limit-collision}
\end{equation}
The convention \(T_0=1\) applies when \(j=2\).
\end{proposition}

\begin{proof}
Choose two real smooth signed tests \(\alpha,\beta\) with zero
integral, supported in disjoint disks, such that
\[
 I=\iint\alpha(u)\log(1/|u-w|)\beta(w)\dd u\dd w\ne0.
\]
Such tests can be obtained by taking horizontal derivatives of
unit mollifiers concentrated about two distinct points on a
horizontal line.  As their radii tend to zero, integration by parts
makes the displayed integral tend to the mixed horizontal derivative
of the logarithmic kernel at these two points, which is nonzero.
Fix one such pair.  For a fixed point \(y_0\), define
\[
 \alpha_r(x)=r^{-2}\alpha((x-y_0)/r),\qquad
 \beta_r(x)=r^{-2}\beta((x-y_0)/r).
\]
Let \(F=\prod_{i=3}^jH(\rho_{z_i,s_i})\), where the spectator disks
are fixed, small, pairwise disjoint and separated from a fixed
closed disk \(\overline{B(y_0,R_0)}\Subset D\).
An empty product is one.  We first prove
\begin{equation}
 \lim_{r\downarrow0}
       \E[H(\alpha_r)H(\beta_r)F]
          =\frac{v}{2\pi}I\,T_{j-2}(\mathbf z).
 \label{eq:limit-plane-factorization}
\end{equation}
We will evaluate this limit in two ways.  A plane copy independent of
the spectators gives the right side of
\eqref{eq:limit-plane-factorization}; the punctured harmonic expansion
will give \(I L(y_0,\mathbf z)\).  Their equality determines \(L\).
For the first evaluation, the plane copy must be independent of the
exterior record that determines the absolute spectator heights.

Fix \(r>0\) small enough that both inner supports lie in
\(B(y_0,Cr)\) and \(\sqrt r<R_0/4\), and return to the finite
mesh along the subsequence defining \(H\).
Use the ambient pair representation of the true law, with both
colors plus on the connected boundary skeleton and height anchored
to zero at one of its sites.  As in Lemma~\ref{lem:boundary-walk},
this gives the required interior height law.
Retain the complementary-spin boundary constraint first as edge
equalities on this connected set.  Conditional on the first spins
and cell statuses, prescribing its single plus component label has
probability exactly \(1/2\).  It therefore creates no additional
bias in a first-color circuit search.

Search from outside for a good first-color circuit in separated
annular slots between scales \(\sqrt r\) and \(R_0/2\).
Lemma~\ref{lem:buffer-circuits} gives failure probability at most
\(Cr^{\kappa_1}\) in the mesh-first limsup, for some
\(\kappa_1>0\).  The exploration always queries a cell status
together with all its first-color sites, using the stopped kernels
of Proposition~\ref{prop:local-mixing}.  Let
\(\mathcal G_\delta\) contain the stopped query record (including its
success or failure outcome), the selected cut and its first-color data
when it exists, and the genuinely exterior pair data needed below.
Do not reveal complementary-spin values at strictly
inside boundary-layer sites merely because their first-color values
were queried.

We verify both the measurability and the conditional law at this
stage.  Each spectator disk can be joined to the zero-height anchor
by paths outside \(B(y_0,R_0)\).  One may choose paths in the ambient
triangulation, outside the search disk, to the anchored skeleton.
Their sums of the local increments from Lemma~\ref{lem:spin-lift}
give the absolute spectator heights, with the correct common lift
fixed at the anchor.  Alternatively, portions of paths in
\(D\setminus\overline{B(y_0,R_0)}\) can detour around the interior
disk and then reach the boundary layer; their fixed compact
portions occur in every sufficiently fine approximation.
The ambient construction requires no microscopic regularity of that
last approach to the skeleton.  Include all these exterior paths
and the pair data on the spectator disks in \(\mathcal G_\delta\).
Thus the entire discrete spectator product \(F_\delta\) is
\(\mathcal G_\delta\)-measurable, as an absolute-height observation,
not merely a spin observation.

Every crossing cell of the selected contour is open with prescribed
first color.  Its factor is constant and leaves the complementary
spin unrestricted.  No imposed equality edge crosses the contour.
By Lemma~\ref{lem:cell-cut}, closed-cell complementary components
on its two sides are disconnected.  Hence revealing exterior labels
reveals no strictly interior component label.
The stopped query convention likewise fixes no further strictly
interior first-color variables.  Conditional on
\(\mathcal G_\delta\), the unexamined interior is exactly the
open-cut law with no inner pins or equalities.
Its height lift may have an exterior-dependent additive constant;
both local tests have exactly zero total and discard that constant.

The cut surrounds a ball of radius comparable to \(\sqrt r\).
The retained-law and exterior-measurability verifications above allow
us to apply Corollary~\ref{cor:conditional-plane-matching}.  It couples
the increments
in \(B(y_0,Cr)\), with its connecting paths, to the fixed plane
increment marginal \(\nu_\delta\) with error
\(C(r/\sqrt r)^{\kappa_2}\), uniformly in each successful exterior
record.  For every such record choose a coupling kernel whose
second marginal is this same \(\nu_\delta\).  On failed searches
choose an independent variable with marginal \(\nu_\delta\).
These are ordinary conditional couplings on finite discrete
spaces; the finite ambient volume may subsequently be exhausted
using the same local comparison.  In particular the resulting
plane variable \(P_\delta\) satisfies
\[
     \Prob(P_\delta\in A\mid\mathcal G_\delta)=\nu_\delta(A)
\]
for every measurable event \(A\), also on failure.  This identity
proves genuine independence of \(P_\delta\) from
\(\mathcal G_\delta\), and therefore from \(F_\delta\).
It is stronger than specifying just its unconditional marginal.

Write \(U_\delta=h_\delta(\alpha_r)\),
\(V_\delta=h_\delta(\beta_r)\), and let \(U'_\delta,V'_\delta\)
be the corresponding statistics of this plane copy.
Matching increments in a connected disk implies equality of both
statistics because their area weights have exact total zero.
If \(E_\delta\) is the union of search failure and coupling mismatch,
then for some \(\kappa>0\)
\[
 \limsup_{\delta\downarrow0}\Prob(E_\delta)\le Cr^\kappa.
\]
Cauchy--Schwarz and equality off \(E_\delta\) give
\begin{align}
 &\left|\E[(U_\delta V_\delta-U'_\delta V'_\delta)F_\delta]\right|
 \nonumber\\
 &\quad\le
 \Prob(E_\delta)^{1/2}
 \left(\|U_\delta V_\delta F_\delta\|_{L^2}
               +\|U'_\delta V'_\delta F_\delta\|_{L^2}\right).
 \label{eq:limit-coupling-error}
\end{align}
Applying H\"older inside each norm at an order at least \(2j\),
\eqref{eq:limit-small-moments} bounds the first norm by
\(C(1+\log(1/r))\).  All fixed spectator norms are bounded.
The plane zero-total moment bound in
Proposition~\ref{prop:plane-moments} bounds the second norm uniformly.
Thus the mesh limsup of \eqref{eq:limit-coupling-error} is at most
\[
                  Cr^{\kappa/2}(1+\log(1/r)),
\]
which tends to zero.

For each fixed \(r\), first let the mesh tend to zero.
Moment convergence treats the true field and the spectators,
whereas independence factors the plane term as
\(\E[U'_\delta V'_\delta]\E[F_\delta]\).
Proposition~\ref{prop:plane-covariance} evaluates its first factor:
\[
 \lim_{\delta\downarrow0}\E[U'_\delta V'_\delta]
   =\frac{v}{2\pi}
       \iint\alpha_r(x)\log(1/|x-y|)\beta_r(y)\dd x\dd y
   =\frac{v}{2\pi}I.
\]
The last equality uses exact zero totals to cancel the dilation
constant.  Separate harmonicity and radial mean value give
\(\E F=T_{j-2}(\mathbf z)\).  Only now let \(r\downarrow0\);
the preceding error bound proves \eqref{eq:limit-plane-factorization}.

We read the same limit from \eqref{eq:limit-puncture}.
First average its spectator variables on their radial disks.
Their harmonic mean value property, including that for \(L\)
from Lemma~\ref{limit:puncture}, evaluates them at \(\mathbf z\).
After substituting \(x=y_0+ru\), \(y=y_0+rw\), the singular term is
\[
 \iint\alpha(u)\beta(w)L(y_0+rw,\mathbf z)
       \{\log(1/r)+\log(1/|u-w|)\}\dd u\dd w .
\]
For every fixed \(w\), the term containing \(\log(1/r)\) is zero
because \(\int\alpha=0\).  This cancellation is exact even though
\(L\) varies with the second colliding variable.
The remaining kernel is bounded on the two disjoint supports;
smoothness of \(L\) makes the remaining integral tend to
\(I L(y_0,\mathbf z)\).
For the regular term, subtract
\(R(y_0;y_0+rw,\mathbf z)\) before integrating against \(\alpha(u)\).
The uniform first-variable derivative bound from
Lemma~\ref{limit:puncture} bounds its contribution by \(Cr\).
The left side of \eqref{eq:limit-plane-factorization} therefore
tends to \(I L(y_0,\mathbf z)\).
Since \(I\ne0\), comparison proves \eqref{eq:limit-collision}.
\end{proof}

\subsection{Wick recursion and convergence of the full law}

The remaining identification uses harmonicity, prescribed collision
singularities, and the boundary condition.  Related collision arguments
identify dimer-height moments in \cite[Section~3]{Kenyon2001} and
six-vertex moments in \cite[Section~7.2]{DCKLM2026}.
Here Proposition~\ref{prop:collision} gives the logarithmic coefficient
from a plane copy independent of the absolute-height spectators, and
Lemma~\ref{lem:no-diagonal} supplies the boundary value and excludes
additional distributions on collisions.

\begin{theorem}[Moment identification and the full field law]
\label{thm:wick}
For every \(j\ge2\) and distinct \(x,y_2,\ldots,y_j\in D\),
\begin{equation}
 T_j(x,y_2,\ldots,y_j)
   =v\sum_{i=2}^j G_D(x,y_i)\,
          T_{j-2}(y_2,\ldots,\widehat{y_i},\ldots,y_j).
 \label{eq:limit-wick}
\end{equation}
Together with \(T_0=1\) and \(T_1=0\), this identifies every moment
distribution with that of \(\sqrt v\,\Phi_D\).
Every subsequential field limit has that law.  Consequently
\(h_\delta/\sqrt v\) converges to \(\Phi_D\) as a random
distribution, with joint convergence of every finite collection of
smooth compactly supported pairings.
\end{theorem}

\begin{proof}
Fix the distinct \(y_i\)'s and subtract the right side of
\eqref{eq:limit-wick} from its left side as a function of \(x\).
It is harmonic outside these finitely many points.
The Dirichlet Green function for the negative Euclidean Laplacian
has local singularity
\[
             \frac1{2\pi}\log(1/|x-y_i|).
\]
Proposition~\ref{prop:collision} shows that its coefficient agrees
with the logarithmic singularity of \(T_j\) at each point.
Lemma~\ref{limit:puncture} therefore makes the difference removable
there.  It is harmonic on all of \(D\).
Equation~\eqref{eq:limit-zero-boundary} gives zero limit at every
boundary point, as does the Dirichlet Green function on a \(C^2\)
Jordan domain.  Thus the difference extends continuously by zero
to the boundary.  The maximum principle, applied to it and to its
negative, makes it identically zero.  This proves
\eqref{eq:limit-wick}.  The base \(T_1=0\) follows from the exact
height-reflection symmetry and uniform integrability.

Induction gives zero for odd moments and, for even \(j\), the sum
over pairings of the products of \(vG_D\).
Lemma~\ref{lem:no-diagonal} says that these are identities of the
entire moment distributions, including tests crossing diagonals.
The Green products are locally integrable; in each pairing their
variables occur in disjoint pairs, and the logarithmic singularity
is integrable in two dimensions.
It follows that for every real \(f\in C_c^\infty(D)\), \(H(f)\)
has exactly the centered Gaussian moments with variance
\[
       v\iint_{D\times D}f(z)G_D(z,w)f(w)\dd z\dd w.
\]
These moments determine its law.  One direct justification avoids
assuming an exponential moment in advance: the nonnegative series
for \(\cosh(tH(f))\), monotone convergence, and the known even
moments give
\(\E\cosh(tH(f))=\exp(t^2\Var(H(f))/2)\) for every real \(t\).
Since \(e^{|tH(f)|}\le2\cosh(tH(f))\), the exponential moment is
finite.  The absolutely integrable power series for its
characteristic function then yields the Gaussian characteristic
function.  This also covers zero variance.

Apply this conclusion to
\(f=\sum_{\ell=1}^k a_\ell f_\ell\), for every real vector
\((a_1,\ldots,a_k)\).  All these real linear combinations are
Gaussian with the stated covariance form.  Their characteristic
functions identify the joint vector
\((H(f_1),\ldots,H(f_k))\) with that of
\(\sqrt v\,\Phi_D\).

Finally choose a countable dense family of smooth tests on an
interior exhaustion of \(D\).  Evaluation on these tests determines
a distribution and generates the Borel sigma-field of the local
Sobolev spaces used in Proposition~\ref{prop:tightness}; this follows,
for example, from their realization by countably many localized
Fourier coordinates.  Hence equality of the finite-dimensional
laws on this family determines the random-distribution law.
We chose the original subsequence arbitrarily, so every limit
allowed by tightness has the same law.  Tightness and uniqueness
give convergence of the full sequence, for each \(t>1\) in
\(H^{-t}_{\mathrm{loc}}(D)\), and therefore as distributions.
Continuity of test evaluation gives the claimed simultaneous
test limits, and Proposition~\ref{prop:tightness} also gives their
mixed-moment convergence.

The normalization is \(\sigma(x)=\sqrt v\).
Its strict positivity, finiteness, and independence of \(D\) and
the approximation were proved in
Proposition~\ref{prop:plane-covariance}.  All estimates here are
for fixed \(x\) in the stated closed interval, so the endpoint is
included with no additional limiting argument in the parameter.
This completes the proof of Theorem~\ref{thm:main}.
\end{proof}
